\documentclass[11pt,a4paper]{article}

% ── Codificación, fuentes y lenguaje ───────────────────────────────────────────
\usepackage[utf8]{inputenc}
\usepackage[T1]{fontenc}
\usepackage[default,scale=0.95]{sourcesanspro}
\renewcommand{\familydefault}{\sfdefault}
\usepackage[spanish,es-noshorthands,es-tabla]{babel}

% ── Geometría y espaciado ─────────────────────────────────────────────────────
\usepackage[top=2.2cm, bottom=2.5cm, left=2.2cm, right=2.2cm, headheight=15pt]{geometry}
\usepackage{setspace}
\setstretch{1.18}
\usepackage{parskip}

% ── Matemáticas y unidades ────────────────────────────────────────────────────
\usepackage{amsmath}
\usepackage{amssymb}
\usepackage{siunitx}

% ── Circuitos ─────────────────────────────────────────────────────────────────
\usepackage[european, straightvoltages]{circuitikz}

% ── Tablas ────────────────────────────────────────────────────────────────────
\usepackage{booktabs}
\usepackage{tabularx}
\usepackage{array}
\usepackage{longtable}
\usepackage{colortbl}
\usepackage{enumitem}

% ── Colores, cajas e iconos ───────────────────────────────────────────────────
\usepackage[dvipsnames]{xcolor}
\usepackage{tcolorbox}
\tcbuselibrary{skins, breakable}
\usepackage{fontawesome5}
\usepackage{graphicx}

% ── Listados de código (dark-mode infográfico) ────────────────────────────────
%   Estilo base: fondo oscuro con números de línea. Los bloques lstlisting se
%   envuelven automáticamente en una caja tcolorbox redondeada.
\usepackage{listings}

% ── Secciones con barra de color (infografía) ────────────────────────────────
\usepackage{titlesec}
\titleformat{\section}
  {\normalfont\LARGE\bfseries\color{azulNoche}}
  {\colorbox{cyanNeon}{\color{fondoOscuro}\thesection}}{0.7em}{}
  [\vspace{2pt}\color{cyanNeon}\rule{\linewidth}{1.8pt}]
\titlespacing*{\section}{0pt}{24pt}{10pt}
\titleformat{\subsection}
  {\normalfont\large\bfseries\color{azulNoche}}
  {\textcolor{cyanNeon}{\thesubsection}}{0.7em}{}
  [\vspace{1pt}\color{grisLinea}\rule{\linewidth}{0.6pt}]
\titlespacing*{\subsection}{0pt}{14pt}{6pt}
\titleformat{\subsubsection}
  {\normalfont\normalsize\bfseries\color{azulNoche}}
  {\textcolor{cyanNeon}{\thesubsubsection}}{0.7em}{}
\titlespacing*{\subsubsection}{0pt}{10pt}{4pt}

% ── Cabeceras y pies de página ────────────────────────────────────────────────
\usepackage{fancyhdr}
\usepackage{lastpage}
\pagestyle{fancy}
\fancyhf{}
\renewcommand{\headrulewidth}{0pt}
\renewcommand{\footrulewidth}{0pt}
\fancyfoot[C]{%
  \begin{minipage}{\linewidth}
    \vspace{2pt}
    \color{cyanNeon}\rule{\linewidth}{1.2pt}\\[2pt]
    \centering\color{grisTexto}\footnotesize
    Página \thepage\ de \pageref{LastPage}
  \end{minipage}%
}

% ── Hiperenlaces ──────────────────────────────────────────────────────────────
\usepackage[colorlinks=true, linkcolor=azulNoche, urlcolor=cyanNeon]{hyperref}
\sloppy
\emergencystretch 3em

% ── Paleta de colores — tecnológico dark-mode ──────────────────────────────────
\definecolor{fondoOscuro}{HTML}{0D1B2A}     % Dark navy — fondo banda título
\definecolor{verdeTurquesa}{HTML}{004D40}    % Verde turquesa — bandas de marca
\definecolor{azulNoche}{HTML}{1B3A6B}       % Midnight blue — secciones, marcos
\definecolor{azulMedio}{HTML}{2A5298}       % Azul medio — variante de acento
\definecolor{cyanNeon}{HTML}{00C8E0}        % Cyan eléctrico — acentos primarios
\definecolor{verdeSignal}{HTML}{00B887}     % Verde señal — fortalezas, éxito
\definecolor{naranjaVivo}{HTML}{FF7043}     % Naranja vivo — mejoras, advertencias
\definecolor{rojoAlerta}{HTML}{E53935}      % Rojo alerta — errores
\definecolor{grisPapel}{HTML}{F4F6FA}       % Fondo tarjetas claras
\definecolor{grisLinea}{HTML}{C8D0E0}       % Bordes sutiles
\definecolor{grisTexto}{HTML}{6B7A99}       % Texto secundario
\definecolor{amarilloNota}{HTML}{FFB300}    % Notas / advertencias doradas

% Colores específicos para bloques de código (dark-mode)
\definecolor{codigoFondo}{HTML}{1E2D3D}      % Fondo del bloque de código
\definecolor{codigoComentario}{HTML}{7A8FA6} % Comentarios
\definecolor{codigoCadena}{HTML}{FF9D5C}     % Cadenas / strings
\definecolor{codigoKeyword}{HTML}{00C8E0}    % Palabras clave
\definecolor{codigoNumero}{HTML}{00E5A0}     % Números

% ── Banda de título: portada infográfica ──────────────────────────────────────
%   Diseño en 2 capas: banda de color (por defecto fondo oscuro) + franja de
%   acento cyan abajo. Uso: \bandaTitulo[color]{título}{subtítulo}.
%   El icono se puede cambiar con \renewcommand{\iconoBanda}{...} (FontAwesome).
\newcommand{\iconoBanda}{microchip}
\newcommand{\bandaTitulo}[3][fondoOscuro]{%
  \begin{tcolorbox}[
    enhanced,
    colback=#1, colframe=#1,
    boxrule=0pt, arc=8pt, outer arc=8pt,
    width=\linewidth,
    top=18pt, bottom=6pt, left=14pt, right=14pt,
    borderline south={3.5pt}{0pt}{cyanNeon},
  ]
    \begin{minipage}[c]{0.07\linewidth}
      \centering
      \textcolor{cyanNeon}{\fontsize{30}{30}\selectfont\faIcon{\iconoBanda}}
    \end{minipage}%
    \hspace{8pt}%
    \begin{minipage}[c]{0.88\linewidth}
      {\fontsize{20}{24}\selectfont\bfseries\color{white}#2}\par\vspace{5pt}%
      \textcolor{cyanNeon!80!white}{\small\faIcon{angle-right}~#3}%
    \end{minipage}
    \vspace{4pt}
  \end{tcolorbox}%
  \vspace{8pt}%
}

% ── Tarjetas de datos (infografía) ────────────────────────────────────────────
\tcbset{
  tarjetaDato/.style={
    enhanced, breakable,
    arc=6pt, outer arc=6pt,
    colback=grisPapel, colframe=azulNoche,
    boxrule=1pt,
    fonttitle=\bfseries\small, coltitle=white,
    attach boxed title to top left={yshift=-3mm, xshift=6mm},
    boxed title style={
      colback=azulNoche, colframe=azulNoche,
      arc=4pt, boxrule=0pt,
    },
    drop fuzzy shadow=azulNoche!30!white,
    top=8pt, bottom=8pt, left=10pt, right=10pt,
  },
  tarjetaIcono/.style={
    enhanced, breakable,
    arc=6pt, outer arc=6pt,
    colback=white, colframe=grisLinea, boxrule=0.8pt,
    fonttitle=\bfseries\small, coltitle=azulNoche,
    borderline west={3pt}{0pt}{cyanNeon},
    drop fuzzy shadow=azulNoche!25!white,
    top=6pt, bottom=6pt, left=10pt, right=10pt,
  },
  cajaTitulo/.style={
    enhanced, breakable,
    arc=6pt, outer arc=6pt,
    colback=azulNoche!10!grisPapel, colframe=azulNoche,
    boxrule=1pt,
    fonttitle=\bfseries, coltitle=white,
    attach boxed title to top left={yshift=-3mm, xshift=6mm},
    boxed title style={
      colback=azulNoche, colframe=azulNoche,
      arc=4pt, boxrule=0pt,
    },
    drop fuzzy shadow=azulNoche!25!white,
    top=8pt, bottom=8pt, left=10pt, right=10pt,
  },
  cajaContenido/.style={
    enhanced, breakable,
    arc=5pt, outer arc=5pt,
    colback=grisPapel, colframe=grisLinea,
    boxrule=0.8pt,
    fonttitle=\bfseries\small, coltitle=azulNoche,
    top=6pt, bottom=6pt, left=10pt, right=10pt,
  },
  cajaFortaleza/.style={
    enhanced, breakable,
    arc=6pt, outer arc=6pt,
    colback=verdeSignal!8!white, colframe=verdeSignal,
    boxrule=1pt,
    borderline west={4pt}{0pt}{verdeSignal},
    fonttitle=\bfseries\small, coltitle=verdeSignal!70!black,
    drop fuzzy shadow=verdeSignal!20!white,
    top=6pt, bottom=6pt, left=10pt, right=10pt,
  },
  cajaMejora/.style={
    enhanced, breakable,
    arc=6pt, outer arc=6pt,
    colback=naranjaVivo!8!white, colframe=naranjaVivo,
    boxrule=1pt,
    borderline west={4pt}{0pt}{naranjaVivo},
    fonttitle=\bfseries\small, coltitle=naranjaVivo!80!black,
    drop fuzzy shadow=naranjaVivo!20!white,
    top=6pt, bottom=6pt, left=10pt, right=10pt,
  },
  cajaRecomendacion/.style={
    enhanced, breakable,
    arc=6pt, outer arc=6pt,
    colback=cyanNeon!6!white, colframe=cyanNeon!60!azulNoche,
    boxrule=1pt,
    borderline west={4pt}{0pt}{cyanNeon},
    fonttitle=\bfseries\small, coltitle=azulNoche,
    drop fuzzy shadow=azulNoche!20!white,
    top=6pt, bottom=6pt, left=10pt, right=10pt,
  },
}

% ── Ícono + texto (encabezados de sección y elementos) ────────────────────────
\newcommand{\iconotexto}[2]{\textcolor{cyanNeon}{\faIcon{#1}}~#2}

% ── Listados de código: estilo dark + auto-envoltura ─────────────────────────
\lstdefinestyle{estiloCodigo}{
    backgroundcolor=\color{codigoFondo},
    basicstyle=\ttfamily\small\color{white},
    commentstyle=\color{codigoComentario}\itshape,
    stringstyle=\color{codigoCadena},
    keywordstyle=\color{codigoKeyword}\bfseries,
    numberstyle=\tiny\color{codigoComentario},
    numbers=left,
    numbersep=10pt,
    showstringspaces=false,
    breaklines=true,
    breakatwhitespace=false,
    tabsize=4,
    frame=none,
    captionpos=b,
    aboveskip=6pt,
    belowskip=4pt,
    xleftmargin=14pt,
    framexleftmargin=14pt,
    literate=
      {á}{{\'a}}1 {é}{{\'e}}1 {í}{{\'i}}1 {ó}{{\'o}}1 {ú}{{\'u}}1
      {Á}{{\'A}}1 {É}{{\'E}}1 {Í}{{\'I}}1 {Ó}{{\'O}}1 {Ú}{{\'U}}1
      {ñ}{{\~n}}1 {Ñ}{{\~N}}1 {ü}{{\"u}}1 {Ü}{{\"U}}1
      {¿}{{?`}}1 {¡}{{!`}}1
}
\lstdefinestyle{estiloPython}{
    style=estiloCodigo,
    language=Python,
}
\lstset{style=estiloCodigo}
\BeforeBeginEnvironment{lstlisting}{%
  \begin{tcolorbox}[
    enhanced, breakable,
    arc=6pt, outer arc=6pt,
    colback=codigoFondo, colframe=azulNoche,
    boxrule=1pt,
    drop fuzzy shadow=azulNoche!25!white,
    top=2pt, bottom=2pt, left=0pt, right=0pt,
  ]%
}
\AfterEndEnvironment{lstlisting}{\end{tcolorbox}}

% ── Cajas de contenido temáticas ──────────────────────────────────────────────
\newtcolorbox{cajaRecuerda}{
    enhanced, breakable,
    arc=6pt, outer arc=6pt,
    colback=azulNoche!10!grisPapel,
    colframe=azulNoche, boxrule=1pt,
    borderline west={4pt}{0pt}{cyanNeon},
    fonttitle=\bfseries\small\color{white},
    title={\faIcon{info-circle}~Recuerda},
    attach boxed title to top left={yshift=-3mm, xshift=6mm},
    boxed title style={colback=azulNoche, arc=4pt, boxrule=0pt},
    drop fuzzy shadow=azulNoche!25!white,
    left=10pt, right=10pt, top=8pt, bottom=8pt,
}

\newtcolorbox{cajaConcepto}{
    enhanced, breakable,
    arc=6pt, outer arc=6pt,
    colback=verdeSignal!8!white,
    colframe=verdeSignal, boxrule=1pt,
    borderline west={4pt}{0pt}{verdeSignal},
    fonttitle=\bfseries\small\color{white},
    title={\faIcon{lightbulb}~Concepto clave},
    attach boxed title to top left={yshift=-3mm, xshift=6mm},
    boxed title style={colback=verdeSignal!80!black, arc=4pt, boxrule=0pt},
    drop fuzzy shadow=verdeSignal!20!white,
    left=10pt, right=10pt, top=8pt, bottom=8pt,
}

\newtcolorbox{cajaEjemplo}{
    enhanced, breakable,
    arc=6pt, outer arc=6pt,
    colback=cyanNeon!5!white,
    colframe=cyanNeon!70!azulNoche, boxrule=1pt,
    borderline west={4pt}{0pt}{cyanNeon},
    fonttitle=\bfseries\small\color{fondoOscuro},
    title={\faIcon{code}~Ejemplo},
    attach boxed title to top left={yshift=-3mm, xshift=6mm},
    boxed title style={colback=cyanNeon, arc=4pt, boxrule=0pt},
    drop fuzzy shadow=azulNoche!20!white,
    left=10pt, right=10pt, top=8pt, bottom=8pt,
}

\newtcolorbox{cajaEjercicio}{
    enhanced, breakable,
    arc=6pt, outer arc=6pt,
    colback=azulMedio!7!white,
    colframe=azulMedio, boxrule=1pt,
    borderline west={4pt}{0pt}{azulMedio},
    fonttitle=\bfseries\small\color{white},
    title={\faIcon{pencil-alt}~Ejercicio},
    attach boxed title to top left={yshift=-3mm, xshift=6mm},
    boxed title style={colback=azulMedio, arc=4pt, boxrule=0pt},
    drop fuzzy shadow=azulNoche!20!white,
    left=10pt, right=10pt, top=8pt, bottom=8pt,
}

\newtcolorbox{cajaReto}{
    enhanced, breakable,
    arc=6pt, outer arc=6pt,
    colback=naranjaVivo!6!white,
    colframe=naranjaVivo, boxrule=1pt,
    borderline west={4pt}{0pt}{naranjaVivo},
    fonttitle=\bfseries\small\color{white},
    title={\faIcon{fire}~Reto},
    attach boxed title to top left={yshift=-3mm, xshift=6mm},
    boxed title style={colback=naranjaVivo, arc=4pt, boxrule=0pt},
    drop fuzzy shadow=naranjaVivo!20!white,
    left=10pt, right=10pt, top=8pt, bottom=8pt,
}

% ── Indicadores de dificultad ─────────────────────────────────────────────────
\newcommand{\facil}{\textcolor{verdeSignal}{\faIcon{circle}\,\textbf{Fácil}}}
\newcommand{\intermedio}{\textcolor{naranjaVivo}{\faIcon{circle}\,\textbf{Intermedio}}}
\newcommand{\dificil}{\textcolor{rojoAlerta}{\faIcon{circle}\,\textbf{Difícil}}}

% ─────────────────────────────────────────────────────────────────────────────

% ============================================================================
% Memo de evaluación: incorporación de características de DeepSeek Harness
% (dsh) a ELECTRONICA — Fase B (exploración y evaluación).
% Generado a partir del skill de consulta `dsh` (docs/SKILL/dsh/dsh.tex)
% y de la especificación docs/AGENTE_IA/dsh-agent-skill-spec.md.
% ============================================================================

% ── Cabecera específica del memo ──
\renewcommand{\iconoBanda}{sitemap}
\fancyhead[L]{\color{azulNoche}\small\bfseries \faIcon{sitemap}~\textcolor{cyanNeon}{dsh $\rightarrow$ ELECTRONICA}}
\fancyhead[R]{\color{grisTexto}\small Evaluación de incorporación -- Fase B}

\begin{document}
\bandaTitulo[fondoOscuro]{DeepSeek Harness en ELECTRONICA}{Evaluación de características del harness dsh y su posible incorporación al espacio de trabajo}

\vspace{6pt}

\begin{tcolorbox}[tarjetaDato, title={\faIcon{clipboard-list}~Resumen ejecutivo}]
\textbf{Objetivo.} Explorar las características de DeepSeek Harness (dsh), compararlas con la
arquitectura determinista de 3 capas de ELECTRONICA y decidir cuáles incorporar.

\textbf{Método.} Se construyó primero el skill de consulta \texttt{dsh}
(vía \texttt{flujo\_repo\_a\_skill.py}, perfil \texttt{referencia\_codigo}) a partir del repo
\texttt{deepseek-ai/deepseek-harness}; este memo se redacta consultando ese skill y la
especificación \texttt{docs/AGENTE\_IA/dsh-agent-skill-spec.md}.

\textbf{Veredicto.} No se recomienda adoptar dsh como segundo harness ni migrar el kernel a
Cordis. Sí se recomienda \textbf{incorporar por analogía} dos ideas concretas: (1) trazabilidad
append-only de sesiones y (2) un gate de validación ``todo es configurable por plugin''.
El detalle está en la matriz de la Sección 3.
\end{tcolorbox}

\section{Qué es DeepSeek Harness}

dsh es un \textbf{agent harness} open-source de DeepSeek AI. Su filosofía central es
\textit{``Everything is a plugin. Every run is traceable.''} y está construido sobre el kernel
\textbf{Cordis} (gestión de plugins, montaje/desmontaje y dependencias).

\begin{itemize}[leftmargin=*,topsep=2pt,itemsep=1pt]
  \item \textbf{Harness vs.\ Modelo}: el modelo es el ``motor de razonamiento''; el harness
        permite al agente interactuar con su entorno, ejecutar herramientas y persistir.
  \item \textbf{Kernel}: Cordis gestiona el ciclo de vida (\texttt{Fiber}), los servicios
        (\texttt{ctx.*}) y los eventos (\texttt{emit/serial/bail/parallel/waterfall}).
  \item \textbf{Trazabilidad}: log de sesión \textbf{append-only} con todos los eventos
        (\texttt{turn/*}, \texttt{step/*}, \texttt{user/message}, \texttt{assistant/message},
        \texttt{tool/*}); habilita reanudar, bifurcar, buscar y reproducir sesiones.
\end{itemize}

\begin{tcolorbox}[cajaTitulo, title={\faIcon{cube}~Los 4 modos de ejecución}]
\textbf{4 modos de ejecución:}
\begin{enumerate}[leftmargin=*]
  \item \textbf{Standard}: agente completo con herramientas locales (edición, shell, búsqueda, subagentes).
  \item \textbf{Code}: tools expuestas vía SDK (TipoScript) para orquestar tool-calls en un programa.
  \item \textbf{Minimal}: benchmark restringido (solo \texttt{bash} y \texttt{str\_replace\_editor}).
  \item \textbf{Creator}: autoría de presets con inspección de runtime y sandboxing de plugins.
\end{enumerate}
\end{tcolorbox}

\section{Comparativa con ELECTRONICA}

\begin{table}[htbp]
\centering
\small
\begin{tabularx}{\linewidth}{>{\raggedright\arraybackslash}p{3.2cm} X X}
\toprule
\textbf{Aspecto} & \textbf{DeepSeek Harness} & \textbf{ELECTRONICA} \\
\midrule
Filosofía & ``Todo es plugin. Cada ejecución es trazable.'' & 3 capas deterministas: Directivas $\rightarrow$ Orquestación $\rightarrow$ Ejecución \\
Kernel & Cordis (DI, ciclo de vida, eventos) & Orquestador opencode + scripts \texttt{execution/*.py} \\
Selección de modelo & Configurable por agente, multi-proveedor & Determinista vía \texttt{enrutador.py} (descriptor $\rightarrow$ tier) \\
Sesiones & Log append-only + resume/fork/search/replay & \texttt{Sessions/*.md} por tema + \texttt{.tmp/run\_state.json} por flujo \\
Extensibilidad & Plugins npm (bundles, parches, overlays) & Scripts deterministas + directivas YAML \\
UI & Web UI local (:3080) & CLI-first (opencode) \\
Lenguaje & TypeScript/pnpm + Cordis + vendor/ & Python (conda) mayoritariamente \\
\bottomrule
\end{tabularx}
\end{table}

\section{Matriz de decisión}

\begin{table}[htbp]
\centering
\small
\begin{tabularx}{\linewidth}{>{\raggedright\arraybackslash}p{3.4cm} X c c}
\toprule
\textbf{Característica dsh} & \textbf{Descripción} & \textbf{Valor} & \textbf{Veredicto} \\
\midrule
Trazabilidad append-only & Log de eventos inmutable por sesión, reanudable & \textbf{Alto} & \textbf{Adaptar} \\
Sistema de plugins / bundles & Toda capacidad es un plugin montable & Medio & Adaptar (parcial) \\
Modos de ejecución & standard/code/minimal/creator & Medio & Descartar (freno) \\
Web UI & Interfaz web local & Bajo & Descarta \\
Framework Cordis como kernel & Sustituir el orquestador por Cordis & Bajo & Descartar \\
Protocolo Typert / remote & RPC tipado Host-Cliente & Bajo & Descartar \\
Gestión de subagentes multi-provider & codex/claude-code/acp/dsh-sdk & Medio & Evaluar a futuro \\
\bottomrule
\end{tabularx}
\end{table}

\subsection{1. Trazabilidad append-only (Alta prioridad -- adaptar)}

dsh garantiza que \textit{todo lo que el modelo ve es reconstruible desde el log}. Hoy
ELECTRONICA guarda el estado de cada flujo multi-paso en \texttt{.tmp/run\_state.json}
(un solo archivo sobrescrito) y las sesiones académicas en \texttt{Sessions/*.md}.

\begin{tcolorbox}[cajaFortaleza, title={\faIcon{check-circle}~Qué incorporar por analogía}]
\begin{itemize}[leftmargin=*,topsep=2pt,itemsep=1pt]
  \item Registrar un \textbf{log de eventos append-only} (\texttt{.tmp/session\_log.jsonl})
        por flujo: \texttt{step/inicio}, \texttt{tool/llamada}, \texttt{resultado}, timestamps.
        Inmutable: nunca se edita ni elimina una entrada.
  \item Derivar \texttt{.tmp/run\_state.json} como vista (checkpoint) de ese log.
  \item Permitir \textbf{reanudar} un flujo interrumpido leyendo el log (resume) y
        \textbf{reproducir} una ejecución (replay) para diagnósticos.
\end{itemize}
\end{tcolorbox}

Impacto: mejora la trazabilidad (principio ``every run is traceable'') sin romper el
determinismo (el log solo registra, no decide). Costo bajo: un helper determinista en
\texttt{execution/}.

\subsection{2. Plugins / bundles (Media prioridad -- adaptar parcialmente)}

La filosofía ``todo es plugin'' de Cordis inspira: cada \textbf{capacidad} (herramienta,
seam, backend) es una unidad montable con configuración y dependencias declaradas.

\begin{tcolorbox}[cajaMejora, title={\faIcon{plug}~Qué tomar prestado}]
\begin{itemize}[leftmargin=*,topsep=2pt,itemsep=1pt]
  \item Formalizar los \textbf{seams} de ELECTRONICA: interfaces abstractas con backends
        intercambiables (LLM, storage, filesystem) detrás de la capa de ejecución.
  \item Los \texttt{execution/*.py} ya son unidades con responsabilidad única: solo falta
        declarar sus dependencias/entradas explícitamente (contrato de CLI) y su ``montaje''
        en la directiva (análogo al \texttt{bundle}).
\end{itemize}
\end{tcolorbox}

Esto \textbf{no implica} adoptar el framework Cordis (TypeScript/npm), sino el \textit{patrón}:
directivas YAML que ``componen'' scripts deterministas.

\subsection{3. Modos de ejecución / Web UI / framework completo (Descartar)}

\begin{cajaRecuerda}
\begin{itemize}[leftmargin=*,topsep=2pt,itemsep=1pt]
  \item Mantener \textbf{dos harness} (opencode + dsh) duplica mantenimiento y consumo de RAM
        (límite estricto de 4 GB del workspace).
  \item La Web UI requiere ejecutar un servicio Node en :3080 de forma persistente; choca con
        el enfoque CLI-first y de scripts deterministas.
  \item Emigrar el orquestador a Cordis rompería la arquitectura de 3 capas validada y todo el
        ecosistema de directivas/scripts ya documentado y en producción.
  \item El proyecto es \textbf{developer preview} con \textit{breaking changes} frecuentes.
\end{itemize}
\end{cajaRecuerda}

\subsection{4. Subagentes multi-provider (Evaluar a futuro)}

La lista \texttt{dsh-subagents} (\texttt{spawn}, \texttt{fork}, \texttt{acp}, \texttt{codex},
\texttt{claude-code}, \texttt{dsh-sdk}) es análoga a nuestra delegación vía \texttt{Task tool} /
subagentes de opencode. No se necesita ahora, pero sí se registra como patrón a considerar si
un día el orquestador delega trabajo a agentes externos de forma sistemática.

\section{Recomendaciones}

\begin{enumerate}[leftmargin=*,topsep=2pt,itemsep=1pt]
  \item \textbf{Corto plazo:} implementar el log append-only como helper determinista en
        \texttt{execution/} y documentar su uso en la directiva genérica de flujos. Ver
        \texttt{directives/} una vez acordado el diseño.
  \item \textbf{Mediano plazo:} formalizar ``seams'' (contratos de interfaz) en unos pocos
        scripts clave (LLM client, storage), sin cambiar su implementación.
  \item \textbf{No hacer ahora:} adoptar dsh/Cordis como kernel, Web UI, o segundo harness.
  \item \textbf{Pendiente del usuario:} prueba hands-on de los 4 modos de dsh
        (\texttt{npx @deepseek-ai/dsh web}) quedó pendiente de decisión.
\end{enumerate}

\begin{tcolorbox}[cajaRecomendacion, title={\faIcon{chess-knight}~Decisión adoptada}]
Incorporar \textbf{solo por analogía} (patrones y principios), no por adopción de la
infraestructura. Se conserva el determinismo como guardrail:
misma entrada $\rightarrow$ misma salida, y la lógica de negocio vive en \texttt{execution/*.py},
nunca en el prompt.
\end{tcolorbox}

\vspace{4pt}
\begin{center}
\color{grisTexto}\small Basado en el skill \texttt{dsh} (2026-09-08) y la especificación
\texttt{dsh-agent-skill-spec.md}. Fase B de la exploración de DeepSeek Harness.
\end{center}

\end{document}